% 中文后封文案（供中文目标使用）

{
\parindent=0pt
\parskip=\baselineskip
{\OPTbacktitlefont
\textit{摘自引言：}}
\OPTbackfont

\emph{同伦类型论}是数学的一个新分支，它以一种出人意料的方式融合了多个不同领域的思想。它建立在近来发现的一条联系之上，即\emph{同伦论}与\emph{类型论}之间的深刻关联。它触及的主题看似彼此遥远，从球面的同伦群，到类型检查算法，再到弱 $\infty$-群胚的定义。

同伦类型论把新的思想引入了数学基础本身。一方面，我们有 Voevodsky 精妙而优美的\emph{泛等公理}。特别地，泛等公理蕴含：同构结构可以被识别。这一原则与传统基础中的“官方”教条并不相容，但数学家在实际工作中早已自然地这样使用。另一方面，我们有\emph{高阶归纳类型}，它为同伦论中的一些基本空间与构造提供了直接的逻辑刻画，例如球面、柱体、截断、局部化等。
这两种思想都无法在经典集合论基础中被直接表达；而在同伦类型论中，它们的结合使一种全新的“同伦类型逻辑”成为可能。

这表明，数学基础可以有一种新的形态：它具有内在的同伦内容，强调数学对象的不变性观点，并且能够被方便地机器实现，从而为工作中的数学家提供切实帮助。这就是\emph{泛等基础}计划。

本书旨在给出泛等基础基本内容的首个系统性阐述，并汇集这种新推理风格的代表性实例。同时，本书不要求读者事先掌握或学习形式逻辑，也不要求使用任何计算机证明助手。
我们相信，泛等基础最终将成为集合论之外的一个可行替代方案，作为大多数数学家在非形式化数学实践中所依赖的“隐式基础”。

\bigskip

\begin{center}
  {\Large
  \textit{可在 HomotopyTypeTheory.org 免费获取本书。}}
\end{center}
}
